\clearpage
\begin{summarybox}
\paragraph{一眼记住：两个零点夹一个“走平的位置”}

\textbf{罗尔定理——两头等高，中间有水平切线：}
\[
\boxed{\begin{gathered}
f(x)\text{ 在 }[a,b]\text{ 连续，在 }(a,b)\text{ 可导},\\
f(a)=f(b)\\[2pt]
\Longrightarrow\quad \exists\,c\in(a,b),\ f'(c)=0.
\end{gathered}}
\]


\textbf{用几何语言读同一条结论：}令 $A(a,h)$、$B(b,h)$，则 $AB$ 水平，且
\[
\boxed{\ k_{AB}=\frac{f(b)-f(a)}{b-a}=0,
\qquad k_{\text{点 }C\text{ 处切线}}=f'(c)=0\ }.
\]
所以\textbf{两点同高 $\Rightarrow$ 中间有与两点连线平行的水平切线}；\textbf{切点不必与两端同高}。

\textbf{零点个数——做题最快的两句：}
\[
\boxed{\begin{aligned}
f\text{ 有 }m\text{ 个不同零点}&\Longrightarrow f'\text{ 至少有 }m-1\text{ 个},\\
f'\text{ 最多有 }k\text{ 个零点}&\Longrightarrow f\text{ 最多有 }k+1\text{ 个}.
\end{aligned}}
\]

\textbf{继续用一次：}如果 $f''(x)$ 在同一区间内始终不等于 $0$，那么 $f(x)$ 在该区间内\textbf{最多有两个不同零点}（前提是 $f$ 在同一区间连续，内部二阶可导）。

\textbf{眼前一亮：把 $x$ 轴换成任何水平线。}若最多只有 $k$ 个水平切线位置，那么对\textbf{任意固定高度 $h$}，方程 $f(x)=h$ 最多只有 $k+1$ 个不同解。\textbf{记忆路线：两点同高 $\rightarrow$ 中间走平 $\rightarrow$ 限制交点。}


\textbf{考试时这样做：}
\begin{enumerate}
  \item \textbf{先看条件：}图像是否连着？区间中间是否处处能求导？
  \item \textbf{求上限：}数 $f'(x)=0$ 有几个不同解；必要时继续求 $f''(x)$。
  \item \textbf{证实存在：}找两处函数值一正一负，让连续图像必须穿过 $x$ 轴。
  \item \textbf{合在一起：}“最多几个”加上“至少几个”，才能得到“恰好几个”。
\end{enumerate}

\textbf{三句防错：}\textbf{不能反推；“最多”不等于“恰好”；不同零点才分别计数。}

\textbf{和导数判断增减的联系：}如果导数始终为正或始终为负，函数本身一直往一个方向走，最多只有一个零点；从罗尔定理也能得到这个结论。
\end{summarybox}


% -------- 新增九种走向的速记 --------
\clearpage
\begin{summarybox}
\paragraph{补充速记：导数为零的九种走向}

\textbf{一点水平，不等于一定掉头。}当 $f'(c)=0$ 时，
如果 $c$ 左、右紧邻两侧的导数分别\textbf{始终为正、始终为零、始终为负}，
就得到九种基本走法：
\[
\begin{array}{c|ccc}
\text{左}\backslash\text{右}&\text{升}&\text{平}&\text{降}\\ \hline
\text{升}&\text{升}\to\text{升}&\text{升}\to\text{平}&\boxed{\text{升}\to\text{降}}\\
\text{平}&\text{平}\to\text{升}&\text{平}\to\text{平}&\text{平}\to\text{降}\\
\text{降}&\boxed{\text{降}\to\text{升}}&\text{降}\to\text{平}&\text{降}\to\text{降}
\end{array}
\]
\textbf{最关键的两类：}先升后降，是最高处；先降后升，是最低处。
而升后仍升、降后仍降，都可能只是在中间\textbf{短暂出现水平切线}。

\textbf{特别说明：}表中的``平''是指一侧\textbf{整小段}保持水平；九种是两侧导数各自保持一种符号时的\textbf{基本组合}，并非无条件穷尽所有函数。

\textbf{把它与罗尔定理分开记：}
\[
\boxed{\text{罗尔定理：两端同高}\ \Rightarrow\ \text{中间至少一处切线水平}}
\]
\[
\boxed{\text{导数为零：只知这一点走平；是否掉头，再看左右}}
\]
\end{summarybox}
